% !TEX root = ../main.tex

\chapter{绪论}

\section{研究背景与意义}

本文选题紧密立足作为能源动力领域核心支撑学科的燃烧科学，所面临的普遍存在且亟待突破的核心基础性挑战——燃烧不稳定性及其诱发的\textbf{熄火}问题。在航空航天、能源发电等关键领域广泛应用的航空发动机、重型燃气轮机等高端动力装备的实际运行过程中，燃烧室的稳定工作边界不仅是决定装备运行安全性、长期可靠性及极限性能充分发挥的关键核心指标，更直接成为制约动力系统推重比进一步提升、燃油经济性持续优化及低排放目标高效实现的核心瓶颈。熄火作为燃烧不稳定性的极端表现形式，往往会引发装备停机、性能骤降甚至结构损坏等严重后果，其本质是流场特性、燃烧反应与外部环境扰动三者之间复杂的强耦合结果，尤其在气流波动、燃料供给扰动等外部激励作用下，火焰稳定裕度会出现显著降低，给高端动力装备的连续安全运行带来难以规避的严重安全隐患，这也使其成为燃烧领域长期以来亟待攻克、备受关注的核心科学问题。

外部扰动中，\textbf{声激励}是工程实际中普遍存在的典型扰动形式。航空发动机、燃气轮机等动力装置在全工况运行区间内，压气机与风扇转子的旋转气动噪声、静叶与动叶的干涉噪声，以及燃烧过程中热释放脉动诱发的热声振荡等，会形成覆盖宽频率范围、幅值动态变化的复杂声学环境。这种声学扰动并非孤立存在，而是与火焰热释放过程形成显著的非线性双向耦合效应，直接调控熄火过程的萌发与扩展。Chaudhuri等人\cite{chaudhuri2012}的研究为该耦合效应提供了关键实验证据，其通过粒子图像测速（PIV）与平面激光诱导荧光（PLIF）技术，揭示了谐波强迫下钝体稳定预混火焰的熄火机制：当上游混合物速度受谐波强迫时，回流区内会产生受迫涡脱落现象，且当施加振荡的对流波长大于流向回流区长度时，涡脱落会直接诱发火焰熄火。

\begin{figure}[!htbp]
	\centering
	
\includegraphics[width=0.8\textwidth]{image1.png}
	\caption{特定扰动频率下的熄火当量比与无扰动时熄火当量比的比值与回流区长度与对流波长比值（$L_{RZ}/\lambda$）的函数关系图\cite{chaudhuri2012}}
	\label{fig:1}
\end{figure}

如图~\ref{fig:1}所示，该图呈现了不同主流速度（$U_m=5, 10, 15$ m/s）下，归一化熄火当量比随回流区平均长度与对流波长比（$L_{RZ}^{mean}/\lambda$，$\lambda=U_m/f$为对流波长，$f$为激励频率）的变化规律。

\begin{equation}
	\varphi(f) = \varnothing(f)/\varnothing_{0}
\end{equation}

其中$\varnothing(f)$为受激励工况熄火当量比，$\varnothing_{0}$为无扰动工况熄火当量比。结果表明，当 $L_{RZ}^{mean}/\lambda$ 处于 0.3~1.0 区间时，归一化熄火当量比达到峰值，且该区间内火焰出现明显的褶皱与断裂现象，证明此比值范围是受迫涡脱落与火焰相互作用最强的区域，直接主导熄火极限的变化。该研究清晰阐释了声激励类谐波强迫对熄火机制的调控作用，也揭示了流动特征参数（回流区长度、对流波长）与熄火特性的定量关联，为本文聚焦声激励下横向射流火焰熄火机理的研究提供了重要理论与实验参考。

\section{国内外研究现状}

\subsection{横向射流火焰流动与燃烧特性}

在工程实践中，\textbf{横向射流火焰}是广泛应用于加力燃烧室、驻涡燃烧室、冲压发动机燃烧室等先进动力装置的核心火焰稳定模式，其本质是燃料射流与主流道内横向高速来流相交融形成的复杂燃烧体系。与简单构型火焰相比，横向射流火焰的流场呈现强湍流、大剪切、多尺度涡结构共存的特征，燃料射流的穿透深度、卷吸混合强度及火焰锚定位置均受主流与射流的动量比、混合特性等因素显著影响。为优化横向射流火焰的混合与稳定性能，本文采用双喷嘴构型（前喷嘴注燃料、后喷嘴注空气），该构型可实现燃料与空气的独立调控，为精准调控混合过程与燃烧模式提供了良好的实验载体。

\begin{figure}[!htbp]
	\centering
	
\includegraphics[width=0.8\textwidth]{image2.png}
	\caption{双喷嘴横向射流火焰的结构示意图\cite{wang2022}}
	\label{fig:2}
\end{figure}

如图~\ref{fig:2}所示，燃料从前喷嘴注入后，先与高温主流初步混合，后与喷嘴注入的空气发生强相互作用，在反向旋转涡对的作用下形成 C 形截面火焰面\cite{wang2022}。这种构型的核心优势在于通过调控喷嘴间距、射流雷诺数及空气/燃料比等参数，可精准调控燃料与空气的混合历程，为研究混合物分数演化与火焰稳定性的关联提供了可控的实验平台。

\begin{figure}[!htbp]
	\centering
	\includegraphics[width=0.8\textwidth]{image3.png}
	\caption{双喷嘴示意图\cite{wang2023}}
	\label{fig:3}
\end{figure}

图~\ref{fig:3}进一步细化了双喷嘴横向射流的混合与自燃过程：燃料先与高温主流混合预热，再与空气耦合，通过延长混合主导的火焰稳定过程实现充分混合，最终在局部当量比与温度满足条件时发生自燃\cite{wang2023}。该示意图清晰呈现了横向射流火焰从燃料注入、混合预热到自燃稳定的完整物理过程，也揭示了混合物分数在空间上的梯度分布特征，这正是影响火焰熄火的关键因素。

\subsection{混合物分数与双色PIV测量技术}

从燃烧本质来看，\textbf{混合物分数}作为基于守恒标量理论的核心参数，依托燃料与氧化剂混合过程中元素守恒的基本特性，不随化学反应进程发生改变，是定量描述局部燃料与空气混合程度及化学当量比分布的关键指标。其空间分布特征与瞬态演化规律直接决定火焰面的精准定位、燃烧热释放率的空间梯度分布及火焰稳定裕度的边界阈值，更是从燃料混合与反应动力学耦合角度解析熄火机理的核心线索。在声激励作用下，横向射流火焰的流场呈现显著的强非定常动态特性：声激励通过周期性的压力脉动，可精准调制流场中大尺度涡结构的生成起始位置、发展演化轨迹与破碎耗散速率，进而改变燃料射流对周围氧化剂的卷吸混合强度，以及混合物分数在时间与空间维度上的动态演化规律。当局部区域的混合物分数因混合失衡偏离可燃极限范围，或混合过程过快导致火焰面拉伸应变率超过临界阈值时，便会率先诱发局部熄火现象；而这些离散的局部熄火区会在非定常流场的扰动下逐步合并、扩展，最终演变为全局性熄火。因此，借助高精度测量手段精准捕捉声激励作用下混合物分数的瞬态全场分布特征，获取其时空演化的定量数据，是揭示声扰动、流场结构、火焰燃烧特性三者耦合致熄机理的核心前提。

然而，在强湍流、非定常声激励耦合的横向射流火焰流场中，混合物分数的测量面临多重技术瓶颈：强湍流导致标量场时空多尺度演化，声激励诱发流场周期性波动与标量输运的强耦合，以及燃烧发光对诊断信号的干扰，均使得传统测量方法难以实现高精度、瞬态全场测量。\textbf{双色示踪粒子图像测速}（双色PIV）技术凭借其非侵入式、高时空分辨率的优势，为解决该瓶颈提供了有效途径：通过在燃料与空气中分别添加不同粒径或不同光学特性的示踪粒子，利用双色滤镜同步捕捉两路粒子的运动信息，可基于守恒标量输运方程反演得到混合物分数的瞬态全场分布，同时获取流场速度信息，实现流场-标量场的同步测量。这种技术能力恰好匹配本文的研究需求，能够为解析声激励下流场扰动与混合物分数演化及火焰应变率变化与熄火的关系完整物理链条提供直接实验证据。

\subsection{横向射流流动结构与混合特性}

在航空航天、能源动力等工程领域，横向射流燃烧技术因具有混合效率高、燃烧强度大等优势，被广泛应用于航空发动机燃烧室、燃气轮机等核心设备中。横向射流火焰的稳定燃烧直接决定了设备的运行效率与安全性，而燃烧不稳定性引发的火焰振荡、熄火等问题，尤其是贫燃条件下的吹熄现象，不仅限制了燃烧系统的贫燃减排潜力，还可能导致燃烧室热损伤甚至设备故障。上游气流的速度脉动作为燃烧系统中普遍存在的非定常扰动形式，其通过改变横向射流的穿透深度、混合特性及火焰结构，直接影响火焰的稳定燃烧状态，是诱发火焰熄火的关键因素之一。

Rock等\cite{rock2020}指出液体燃料的近贫燃熄火伴随熄灭-复燃循环，高DCN燃料更难熄火且复燃率更高，多数复燃依赖热产物回流；Wan和Zhao\cite{wan2020}发现雷诺数决定支架稳定火焰的熄火模式，大雷诺数下会出现火焰夹断，小雷诺数则火焰因根部下游迁移、热量流失而熄灭；Yi等\cite{yi2019}揭示横流射流中非预混火焰的熄火与自燃核演化相关，迎风侧因高应变率熄火位置更靠后；Yin等\cite{yin2017}表明受限湍流火焰的近贫燃熄火由自持射流振荡引发，慢节奏振荡更易导致熄火-复燃循环，整体可见熄火行为普遍受燃料特性、流动参数和火焰结构的耦合影响。

横向射流的流动结构由射流与主流的相互作用主导，其核心特征包括射流轨迹偏转、剪切层涡旋演化及回流区形成，这些流动特性直接决定了燃料与氧化剂的混合效率，进而影响燃烧过程的稳定性。早期研究中，Fric和Roshko\cite{fric1994}通过流动可视化实验，首次明确了横向射流尾迹中存在反向旋转涡对，该结构通过卷吸主流流体促进射流扩散，是调控混合过程的关键因素。后续Muppidi和Mahesh\cite{muppidi2005}基于直接数值模拟方法，建立了射流轨迹的预测模型，量化分析了射流动量比和密度比对轨迹偏转程度的影响规律，为横向射流流动特性的定量描述提供了理论基础。

Zhang和Yang\cite{zhang2017a,zhang2017b}通过实验与数值模拟相结合的方式，设置射流参数（射流中心速度160~m/s，横流速度40~m/s），工作介质为空气，射流直径 $d=1.27$~mm。分为a（横流无激励）、b（激励频率5~kHz，速度振荡振幅10\%）、c（激励频率5~kHz，速度振荡振幅50\%）三个对比案例，得到如图~\ref{fig:4}所示的结果。图~\ref{fig:4}展示了射流中心平面的标量浓度分布，核心反映射流与横流的混合效率、射流扩散范围及浓度均匀性，是标量输运过程的直接体现。

\begin{figure}[!htbp]
	\centering
	\includegraphics[width=0.9\textwidth]{image4.png}
	\caption{射流中心平面标量浓度的时间演化\cite{zhang2017b}}
	\label{fig:4}
\end{figure}

系统研究了横向射流近场的流动动力学与混合特性，发现射流出口处的剪切层卷起与主流剪切作用共同主导了初始涡旋的形成，而涡旋的破碎与融合过程直接影响混合的均匀性，为理解横向射流的混合机制提供了重要参考。Gevorkyan等\cite{gevorkyan2016}通过浓度测量实验进一步证实，横向射流的混合过程具有明显的区域差异性，近场混合主要由涡旋卷吸作用主导，远场则以湍流扩散为主，且混合效率随射流穿透深度的增加而显著提升。

\subsection{剪切层不稳定性与流动控制}

剪切层不稳定性是横向射流流动演化的核心驱动力，相关研究已形成系列成果。Megerian等\cite{megerian2007}通过线性稳定性分析，识别出横向射流剪切层中存在Kelvin-Helmholtz（K-H）不稳定性和Orr-Sommerfeld（O-S）不稳定性两种主导模式，其中K-H不稳定性主导剪切层的初始卷起，O-S不稳定性则影响涡旋的后续发展与破碎过程。DeB.Alves等\cite{alves2008}在此基础上通过实验验证，发现剪切层不稳定性的特征频率与射流出口速度和主流速度的比值密切相关，为通过调控流动参数抑制不稳定性提供了关键依据。针对流动控制需求，Harris等\cite{harris2021}研究了tabs结构对横向射流不稳定性的调控作用，结果表明合理设计的tabs能够抑制特定频率的不稳定性，延缓涡旋破碎，从而改善射流混合的均匀性与稳定性。此外，Getsinger等\cite{getsinger2014}通过PIV技术研究了受限空间内横向射流的结构与稳定性特性，发现边界约束会增强射流与主流的相互作用，导致涡旋结构更复杂，混合效率显著提升，为复杂边界条件下横向射流的应用提供了参考。

横向射流火焰的燃烧稳定性与流动特性紧密耦合，旋流的引入进一步增强了这种耦合效应，同时也增加了燃烧不稳定性的复杂性。Candel等\cite{candel2012}在旋流火焰动力学的综述研究中指出，旋流诱导的回流区能够稳定火焰根部，为燃烧提供持续的热量与自由基，但回流区的动态演化也可能引发燃烧不稳定性，表现为压力振荡与热释放率振荡的相互耦合。Candel等\cite{candel2014}进一步深入研究了旋流火焰的动力学特性，发现当旋流数超过临界值时，回流区会发生涡破裂现象，导致火焰根部位置偏移，进而诱发强烈的燃烧振荡。Huang和Yang\cite{huang2009}针对贫燃预混旋流稳定火焰，系统分析了流动不稳定性对火焰结构的影响，证实了涡破裂引发的回流区形态变化是导致火焰不稳定的关键因素，为旋流火焰的稳定性控制提供了理论支撑。

Palies等\cite{palies2010}通过实验与数值模拟结合，揭示了旋流器动态特性与湍流预混火焰的耦合机制，认为旋流器的非定常流动会直接导致热释放率波动，进而诱发燃烧不稳定性，该研究为旋流器结构的优化设计提供了重要参考。燃烧不稳定性的本质是热释放率振荡与声学振荡的正反馈循环，Lieuwen\cite{lieuwen2003}建立了预混燃烧-声波相互作用的理论模型，量化了热释放率对声学扰动的响应特性，为燃烧不稳定性的预测与抑制提供了核心理论框架。Kim等\cite{kim2008}通过实验研究了火焰结构对火焰传递函数的影响，发现火焰面的拉伸和弯曲会改变热释放率的空间分布，进而影响FTF的幅值和相位，为燃烧不稳定性的定量分析提供了关键参数支持。针对液体燃料燃烧，Sim等\cite{sim2020}研究了声学强制下液滴燃烧的动态响应，发现声学扰动会改变液滴的蒸发速率和燃烧强度，进而影响整体热释放率，该结果为多相燃烧不稳定性的研究提供了重要参考。

\subsection{熄火机理与上游速度脉动}

熄火机理方面，贫燃吹熄（LBO）是横向射流火焰最常见的熄火形式，其核心是局部燃料/空气混合比偏离可燃范围或火焰温度低于临界燃烧温度，导致燃烧反应无法持续进行。Shanbhogue等\cite{shanbhogue2009}通过实验建立了钝体稳定火焰的LBO标度律，发现LBO极限与燃料的化学计量比、流动雷诺数及回流区长度密切相关，为LBO极限的预测提供了定量方法。Broadwell和Breidenthal\cite{broadwell1982}提出的简单混合与化学反应模型，揭示了混合时间与反应时间的匹配关系对燃烧完整性的影响，当混合时间远大于反应时间时，易发生局部熄火甚至整体吹熄，该模型为熄火机理的分析提供了基础理论支撑。Peters\cite{peters1988}提出的层流小火焰概念，进一步完善了湍流燃烧中混合与反应耦合的理论体系，明确了湍流脉动对火焰结构及熄灭过程的影响机制，为熄火边界的精准预测提供了重要理论依据。

上游速度脉动作为燃烧室中的典型非定常扰动，其对横向射流火焰熄火的影响逐渐受到关注。Karagozian等\cite{karagozian2003}系统研究了横向射流的主动控制策略，发现上游流动扰动可通过改变射流的穿透深度和混合特性，间接影响火焰稳定性，该研究为通过调控上游流动参数改善火焰稳定性提供了新思路。Karagozian\cite{karagozian2006}在混合与反应流控制的综述中指出，上游扰动的频率和幅值是影响横向射流混合及燃烧稳定性的关键参数，但相关定量研究仍需进一步深入。Karagozian\cite{karagozian2010}在横向射流及其控制的研究中，进一步证实了上游流动调控对火焰稳定性的显著影响，为上游速度脉动相关研究奠定了重要基础。Karagozian\cite{karagozian2014}在横向射流的综述研究中，通过烟雾可视化技术，直观对比了无激励、正弦波激励和方波激励三种条件下，横流中射流的流动形态差异，核心目的是验证不同激励方式对射流穿透深度、扩散范围及涡结构形成的影响，为后续基于流动稳定性的射流控制策略提供实验依据，其实验结果如图~\ref{fig:5}所示。强调了上游非定常扰动对射流火焰稳定性影响的研究缺口，指出需深入揭示其内在作用机理。孙等\cite{sun2023}在航空发动机燃烧不稳定性研究进展中明确指出，上游气流的非定常扰动是引发燃烧振荡和熄火的重要诱因，但目前关于其定量影响规律和内在机理的研究仍不充分，亟待进一步深入探索。

\begin{figure}[!htbp]
	\centering
	\includegraphics[width=0.9\textwidth]{image5.png}
	\caption{无激励、正弦波、方波对横向射流影响\cite{karagozian2014}}
	\label{fig:5}
\end{figure}

Allison等\cite{allison2015}率先采用PLIF技术突破传统测量局限，为火焰几何与不稳定性的关联研究提供了精准工具；在此基础上，Deng等\cite{deng2019}聚焦声学参数这一关键外部因素，揭示了其对燃烧污染物排放的调控规律，而Liu等\cite{liu2023}则进一步深入不稳定性的触发机制，明确了低频振荡的阈值特性与耦合过程。针对更复杂的燃烧系统，Bonciolini和Noiray\cite{bonciolini2019}探讨了顺序燃烧器中热声模态的同步现象，凸显了熵波这一特殊驱动因素的作用。

\subsection{混合增强与PIV诊断技术}

混合增强技术（尤其是通过射流调制实现的混合增强）一直是研究热点。Vermeulen等\cite{vermeulen1990}采用扬声器驱动器，在宽激励振幅和射流雷诺数范围内，对受限横流中的射流进行激励，结果在喷射孔处产生了一系列环形涡旋，显著提高了射流的扩散范围与穿透深度，并发现当斯特劳哈尔数为0.22时，混合响应达到最优。

斯特拉哈尔数定义为：
\begin{equation}
	S_t = \frac{fd}{U_j}
\end{equation}
其中$f$为激励频率，$d$为射流直径，$U_j$为射流速度。Gogineni等\cite{gogineni1998}通过在管道侧壁布置四个压电致动器对矩形射流进行激励，使低速尾迹区域收缩，射流穿透深度增加。Johari等\cite{johari1999}采用电磁阀在喷射周期的部分时段切断水流供应，发现射流的完全调制会不同程度地改变射流结构、穿透深度、稀释度和混合效果，具体依赖于脉冲持续时间与频率。Eroglu与Breidenthal\cite{eroglu2001}同样通过电磁阀产生方波激励，得出结论：在近场区域可实现最优的涡环间距与强度，从而提升射流穿透深度与混合效率。M'Closkey等\cite{mcloskey2002}量化了时间强制圆形射流的驱动动力学特性，并开发了一种用于开环射流控制的射流补偿系统设计方法，观察到当方波激励频率为射流自然涡脱落频率的次谐波时，射流的穿透深度与扩散范围达到最优。方波激励是提升射流穿透和扩散的核心，正弦波激励难以形成强穿透的离散涡结构，混合增强效果微弱。其实验工况如表~\ref{tab:1}所示，实验结果如图~\ref{fig:6}所示。Narayanan等\cite{narayanan2003}采用开环控制，在射流斯特劳哈尔数为0.1的条件下进行激励，发现该控制方法能有效规整非定常性并增强混合。Shapiro等\cite{shapiro2006}针对多种激励频率，确定了方波激励的最优时间脉冲宽度。

\begin{table}[!htbp]
	\centering
	\caption{实验工况表\cite{mcloskey2002}}
	\label{tab:1}
	\begin{tabular}{c|c|c|c|c|c|c|c|c|c}
		\hline
		编号 & a & b & c & d & e & f & g & h & i \\
		\hline
		激励类型 & 无 & 正弦 & 正弦 & 方波 & 方波 & 方波 & 方波 & 方波 & 方波 \\
		激励参数 & 无 & 73.5Hz & 73.5Hz & 110Hz & 110Hz & 55Hz & 73.5Hz & 85Hz & 220Hz \\
		\hline
	\end{tabular}
\end{table}

\begin{figure}[!htbp]
	\centering
	\includegraphics[width=0.8\textwidth]{image6.png}
	\caption{实验结果图\cite{mcloskey2002}}
	\label{fig:6}
\end{figure}

粒子图像测速（PIV）技术作为全场流动测量的核心手段，凭借其非侵入性、高分辨率等优势，已广泛应用于横向射流和燃烧流动的研究中，为流动特性的精准表征提供了关键技术支撑。Sakakibara和Adrian\cite{sakakibara1999}基于PIV技术实现了水相中温度的全场测量，通过示踪粒子的响应特性与温度的关联，建立了温度测量的定量方法，为后续燃烧流场中温度和速度的同步测量提供了重要技术参考。Kähler等\cite{kahler2016}在第4届国际PIV挑战的总结中，系统评估了不同PIV测量方案的精度和适用性，针对复杂燃烧流场的测量难点，提出了优化策略，为PIV技术在燃烧领域的应用提供了方法学指导。

Kashanj和Nobes\cite{kashanj2023}采用4D双色LIF技术研究侧向受限瑞利-贝纳德对流，优化了温度测量灵敏度，强调体积测量对准确分析传热特性的重要性。King等\cite{king1997}利用双示踪剂PLIF技术，实现了湍流剪切层中分子混合程度的瞬时定量检测。Kosiwczuk等\cite{kosiwczuk2005}开发了基于荧光成像的两相PIV技术，成功实现气液两相速度场的同时独立测量。Ulrich等\cite{ulrich2024}将双色LIF技术应用于水/乙醇混合液滴，兼顾了温度测量与成分分析，拓展了技术的工业应用场景。这些研究均围绕LIF、PLIF或PIV相关光学诊断技术展开，分别适配单相对流、湍流混合、两相流及混合液滴等不同研究场景，验证了此类技术在流场参数测量中的核心价值，为相关领域的后续研究提供了技术参考。

Brübach等\cite{brubach2013}测得多种磷光体的温度-寿命特性，其即热成像磷光体受激后，发射光的衰减寿命随温度变化的规律，其结果如图~\ref{fig:7}所示。揭示不同材料在宽温域（近绝对零度至1970~K）的响应差异。结合双色示踪测速，可通过优选磷光体，利用其温度敏感发光特性，关联上游速度脉动与横向射流火焰温度场，助力熄火机理研究。

\begin{figure}[!htbp]
	\centering
	\includegraphics[width=0.9\textwidth]{image7.png}
	\caption{多种热释光荧光粉的温度寿命特性\cite{brubach2013}}
	\label{fig:7}
\end{figure}

Aldén等\cite{alden2011}综述了热成像磷光体在燃烧等场景的测温应用，展示了多种磷光体的温度敏感性，结果如图~\ref{fig:8}所示，实线为寿命-温度关系，虚线为光谱强度比-温度关系，曲线斜率反映灵敏度。结合研究得到结论，热成像磷光体可作为双色示踪粒子，其温度敏感性可助力同步获取横向射流火焰温度与速度场。适配的磷光体能在火焰环境下稳定响应，为解析速度脉动与温度场耦合对熄火的影响提供关键数据支撑，提升熄火机理研究的准确性。

\begin{figure}[!htbp]
	\centering
	\includegraphics[width=0.8\textwidth]{image8.png}
	\caption{多种荧光粉的温度敏感性\cite{alden2011}}
	\label{fig:8}
\end{figure}

Abram等\cite{abram2018}使用BAM:Eu$^{2+}$磷光粉末为样本，采用脉冲Nd:YAG激光的355~nm波长激发，通过光谱仪与探测器采集其室温下的激发和发射光谱，其结果如图~\ref{fig:9}所示。激发光谱在 355~nm 附近呈现强吸收峰，匹配常用激光光源；发射光谱为中心 445~nm 的宽谱带。实验中同步记录激发峰与发射峰的强度及波长分布，明确其光谱响应特性：宽激发谱适配多种激光，可见区强发射便于探测器捕捉，且光谱稳定性不受流体环境影响。该表征为双色成像实验中滤镜的选择提供依据，通过分割其发射谱带获取强度比，进而实现温度测量，适配横向射流火焰的复杂工况，为速度脉动与熄火关联分析提供可靠的温度示踪基础。

\begin{figure}[!htbp]
	\centering
	\includegraphics[width=0.8\textwidth]{image9.jpg}
	\caption{BAM粉末室温激发与发射光谱\cite{abram2018}}
	\label{fig:9}
\end{figure}

双色示踪PIV技术作为PIV的拓展形式，通过两种不同响应特性的示踪粒子（如不同材质、不同粒径或不同荧光特性的粒子），能够实现速度与温度的同步测量，有效解决了单一测量技术难以捕捉多物理量耦合效应的问题。Oberleithner等\cite{oberleithner2012}利用双色示踪技术研究了湍流涡破裂的间歇性特征，通过两种示踪粒子的速度差异和荧光强度变化，同步获取了速度场和温度场信息，展示了该技术在非定常流动测量中的独特优势。该技术能够精准捕捉流动与燃烧的耦合演化过程，为上游速度脉动下横向射流火焰熄火机理的研究提供了理想的测量手段。

Polifke等\cite{polifke2002}通过PIV技术研究了预混燃烧器的不稳定性，发现非单调压力响应是燃烧不稳定性的重要诱因，该研究证实了PIV技术在燃烧不稳定性机理分析中的有效性。Shapiro等\cite{shapiro2006}在受控横向射流的优化研究中，利用PIV技术精准测量了不同控制策略下的流场结构变化，为流动控制效果的评估提供了定量数据支撑。这些研究充分验证了PIV技术在燃烧流场研究中的可靠性，为本文采用双色示踪PIV技术开展相关研究提供了坚实的技术基础。

\section{本文研究内容}

本文以双喷嘴横向射流火焰为研究对象，以上游速度脉动为扰动输入，以声激励为主动调控手段，系统探究上游速度脉动对横向射流火焰熄火边界的影响机理。主要研究内容包括：

（1）搭建集成声激励系统的两级串联横向射流燃烧实验平台，设计加工四组不同孔径与间距的双喷嘴，建立基于双色PIV的流场-混合物分数同步测量方法；

（2）完成声激励脉动系统的设计、加工与调试，实现宽频范围内正弦波与方波激励的稳定输出，并筛选一级燃烧室稳定工况，为二级实验提供可靠热来流；

（3）开展双色PIV光路搭建与数据处理流程开发，实现红蓝粒子的自适应分离、速度场求解与混合物分数计算；

（4）通过前置实验明确喷嘴孔径对燃烧的影响，确立小孔径双喷嘴正式实验方案；

（5）采用DMD等方法对冷态流场进行分析，识别主导模态，为后续声激励实验提供基准；

（6）系统开展无声激励和有声激励条件下的熄火边界测量，揭示声激励参数对熄火极限的调控规律，并从流场-混合协同诊断角度建立“速度脉动→涡结构演变→混合物分数变化→火焰响应”的物理分析框架。